\section{Complete Pipeline and Implementation}

By now we have covered every component of the feature-based face morphing algorithm. In this chapter we bring everything together: a complete worked-out pseudocode specification, a full Python implementation supporting both \texttt{dlib} and \texttt{MediaPipe} backends for landmark detection, a complexity analysis, and practical debugging tips.

\subsection{Full Algorithm Specification}

Algorithm~\ref{alg:face-morph} presents the complete face morphing procedure. The algorithm takes two input images $I_A$ and $I_B$, a choice of blending parameter $\alpha \in [0,1]$, and produces the morphed output image $O$.

\begin{algorithm}[t]
\caption{Feature-Based Face Morphing}
\label{alg:face-morph}
\begin{algorithmic}[1]
\Require Two face images $I_A, I_B$; blending weight $\alpha \in [0,1]$
\Ensure Morphed image $O$

\State $L_A \gets \texttt{detectLandmarks}(I_A)$ \Comment{68-point landmarks}
\State $L_B \gets \texttt{detectLandmarks}(I_B)$
\State $M \gets \{(a_i + b_i)/2 \mid (a_i, b_i) \in (L_A, L_B)\}$ \Comment{Mean shape}
\State $\texttt{triangles} \gets \texttt{delaunayTriangulate}(M)$
\State $O \gets \texttt{zeros like } I_A$
\ForAll{triangle $(i, j, k)$ in \texttt{triangles}}
    \State $T_M \gets (M_i, M_j, M_k)$ \Comment{Triangle in mean shape}
    \State $T_A \gets (L_{A,i}, L_{A,j}, L_{A,k})$
    \State $T_B \gets (L_{B,i}, L_{B,j}, L_{B,k})$
    \State $H_A \gets \texttt{computeAffineTransform}(T_A, T_M)$
    \State $H_B \gets \texttt{computeAffineTransform}(T_B, T_M)$
    \State Compute $H_A^{-1}, H_B^{-1}$ \Comment{Inverse warp matrices}
    \State $[\texttt{xmin}, \texttt{xmax}, \texttt{ymin}, \texttt{ymax}] \gets \texttt{boundingBox}(T_M)$
    \For{$y$ from $\texttt{ymin}$ to $\texttt{ymax}$}
        \For{$x$ from $\texttt{xmin}$ to $\texttt{xmax}$}
            \State \textcolor{gray}{\texttt{// Inverse map to source coordinates}}
            \State $(x_A, y_A) \gets H_A^{-1} \cdot (x, y, 1)$
            \State $(x_B, y_B) \gets H_B^{-1} \cdot (x, y, 1)$
            \State \textcolor{gray}{\texttt{// Sample with bilinear interpolation}}
            \State $c_A \gets \texttt{bilinearSample}(I_A, x_A, y_A)$
            \State $c_B \gets \texttt{bilinearSample}(I_B, x_B, y_B)$
            \State \textcolor{gray}{\texttt{// Alpha blend}}
            \State $O[x,y] \gets \alpha \cdot c_A + (1-\alpha) \cdot c_B$
        \EndFor
    \EndFor
\EndFor
\State \texttt{// Optional: fill background (outside convex hull)}
\State $O \gets \texttt{fillBackground}(O, I_A, I_B, M, L_A, L_B)$
\State \Return $O$
\end{algorithmic}
\end{algorithm}

The algorithm has a clear structure: detect landmarks (lines~1--2), compute the mean shape (line~3), build the Delaunay triangulation (line~4), then iterate over each triangle to warp and blend (lines~5--22). The inverse mapping at lines~14--15 ensures that every output pixel receives a value, avoiding the hole problem discussed in Section~\ref{sec:inverse-mapping}.

\subsection{Morph Sequence Generation}

Algorithm~\ref{alg:face-morph} produces a single frame for a fixed $\alpha$. To generate a complete morph animation from Face A to Face B, we wrap the algorithm in a frame generation loop:

\begin{algorithm}[t]
\caption{Generate Morph Sequence}
\label{alg:morph-sequence}
\begin{algorithmic}[1]
\Require Images $I_A, I_B$; number of frames $N$
\For{$k$ from $0$ to $N-1$}
    \State $t \gets k / (N - 1)$ \Comment{Interpolation parameter}
    \State $M_t \gets \{(1-t) \cdot a_i + t \cdot b_i\}$ \Comment{Interpolated shape}
    \State $\texttt{triangles} \gets \texttt{delaunayTriangulate}(M_t)$
    \State $O_t \gets \texttt{piecewiseAffineWarp}(I_A, I_B, M_t, t)$
    \State \texttt{saveFrame}$(O_t, k)$
\EndFor
\end{algorithmic}
\end{algorithm}

Note that for a morph sequence, the triangulation is recomputed at each frame on the interpolated shape $M_t$, not on the mean shape $M_{0.5}$. This ensures that the triangle structure adapts smoothly as the face geometry transitions.

\subsection{Complete Python Implementation}

We now present the full Python implementation. The code supports two landmark backends:

\begin{itemize}
    \item \textbf{dlib}: The classic 68-point detector. Fast, well-tested, requires the shape predictor model file.
    \item \textbf{MediaPipe}: Google's 478-point 3D mesh. Denser landmarks, real-time performance, but landmarks must be mapped to 68-point indices for compatibility with our Delaunay approach.
\end{itemize}

Both backends share the same core warping and blending code (Algorithms~\ref{alg:face-morph} and~\ref{alg:morph-sequence}). The landmark backend is selected via a command-line argument.

\subsubsection*{Imports and Setup}

\begin{minted}{python}
#!/usr/bin/env python3
"""Feature-based face morphing with dlib and MediaPipe backends."""

import argparse
import os
import cv2
import numpy as np
from scipy.spatial import Delaunay
\end{minted}

\subsubsection*{Landmark Detection}

\paragraph{MediaPipe backend.}

\begin{minted}{python}
def get_landmarks_mediapipe(image):
    """Detect 68-point landmarks using MediaPipe Face Mesh."""
    import mediapipe as mp
    mp_face_mesh = mp.solutions.face_mesh
    # MediaPipe uses 478 3D landmarks. We use a mapping to select 68 points.
    # This mapping is based on the correspondence between MediaPipe indices
    # and the iBUG 68-point model.
    MEDIAPIPE_TO_68 = [
        162, 234, 93, 58, 172, 136, 149, 148, 152, 377, 378, 365,
        397, 288, 323, 454, 389,
        70, 63, 105, 66, 107,  # left eyebrow
        336, 296, 334, 293, 300,  # right eyebrow
        168, 6, 197, 195, 5,  # nose bridge
        48, 49, 98, 97, 2, 326, 327, 309,  # nose base
        33, 160, 158, 133, 153, 144,  # left eye
        362, 385, 387, 263, 373, 380,  # right eye
        61, 146, 91, 181, 84, 17, 314, 405, 321, 375,  # outer mouth
        78, 95, 88, 178, 87, 14, 317, 402, 318, 324,  # inner mouth
        199  # chin center (optional)
    ]

    h, w = image.shape[:2]
    with mp_face_mesh.FaceMesh(
            static_image_mode=True,
            max_num_faces=1,
            refine_landmarks=True) as face_mesh:
        results = face_mesh.process(cv2.cvtColor(image, cv2.COLOR_BGR2RGB))

    if not results.multi_face_landmarks:
        raise ValueError("No face detected in image")

    landmarks = results.multi_face_landmarks[0].landmark
    points = []
    for idx in MEDIAPIPE_TO_68:
        pt = landmarks[idx]
        points.append((pt.x * w, pt.y * h))
    return np.array(points, dtype=np.float64)
\end{minted}

\paragraph{dlib backend.}

\begin{minted}{python}
def get_landmarks_dlib(image, predictor_path):
    """Detect 68-point landmarks using dlib."""
    import dlib

    detector = dlib.get_frontal_face_detector()
    predictor = dlib.shape_predictor(predictor_path)

    gray = cv2.cvtColor(image, cv2.COLOR_BGR2GRAY)
    faces = detector(gray, 1)
    if len(faces) == 0:
        raise ValueError("No face detected in image")

    shape = predictor(gray, faces[0])
    points = np.array(
        [[shape.part(i).x, shape.part(i).y] for i in range(68)],
        dtype=np.float64
    )
    return points
\end{minted}

\subsubsection*{Core Geometry: Affine Transform}

\paragraph{Computing the affine warp.}
Given three source points and three destination points, compute the $2 \times 3$ affine transformation matrix:

\begin{minted}{python}
def compute_affine_transform(src_pts, dst_pts):
    """Compute 2x3 affine matrix mapping src_pts -> dst_pts.

    Both src_pts and dst_pts are 3x2 numpy arrays.
    Returns a 2x3 matrix.
    """
    src_h = np.vstack([src_pts.T, [1, 1, 1]])  # 3x3 homogeneous
    dst_h = np.vstack([dst_pts.T, [1, 1, 1]])
    T = dst_h @ np.linalg.inv(src_h)  # 3x3
    return T[:2, :].astype(np.float32)  # 2x3
\end{minted}

\paragraph{Computing the inverse affine warp.}
For inverse mapping we need $H^{-1}$, the inverse of the affine matrix:

\begin{minted}{python}
def compute_inverse_warp(src_pts, dst_pts):
    """Compute the inverse warp matrix: for each pixel in dst,
    find the corresponding source coordinate."""
    return compute_affine_transform(dst_pts, src_pts)
\end{minted}

\subsubsection*{Delaunay Triangulation}

\begin{minted}{python}
def get_delaunay_triangles(points):
    """Compute Delaunay triangulation and return triangle indices."""
    tri = Delaunay(points)
    return tri.simplices  # Nx3 array, each row: 3 point indices
\end{minted}

For morph sequences where the geometry changes every frame, this function is called on each interpolated shape. For a static average ($\alpha=0.5$), it is called once on the mean shape.

\subsubsection*{Piecewise Affine Warping and Blending}

This is the heart of the algorithm. For each triangle, we build the affine warp from each source face to the mean shape, then apply inverse mapping with bilinear interpolation.

\begin{minted}{python}
def warp_triangle(src_img, src_tri, dst_tri, size):
    """Warp a single triangle from src_img to destination.

    src_tri: 3x2 array of source triangle vertices
    dst_tri: 3x2 array of destination triangle vertices
    size: (width, height) of the output image

    Returns: a warped image (same size as output, zeros outside this triangle)
    """
    mask = np.zeros((size[1], size[0]), dtype=np.uint8)
    cv2.fillConvexPoly(mask, dst_tri.astype(np.int32), 255)

    # Compute inverse warp: dst -> src
    warp_mat = compute_inverse_warp(src_tri, dst_tri)

    # Apply affine warp (OpenCV handles bilinear interpolation internally)
    warped = cv2.warpAffine(
        src_img, warp_mat, (size[0], size[1]),
        borderMode=cv2.BORDER_REFLECT_101
    )

    # Mask to this triangle only
    result = np.zeros_like(src_img)
    for c in range(src_img.shape[2]):
        result[:, :, c] = warped[:, :, c] * (mask > 0)

    return result
\end{minted}

The main blending loop:

\begin{minted}{python}
def morph_faces(img1, img2, landmarks1, landmarks2, alpha=0.5):
    """Morph two faces. alpha=0.5 produces the average face."""
    # Compute interpolated shape
    pts1 = landmarks1.astype(np.float64)
    pts2 = landmarks2.astype(np.float64)
    pts_mean = (1 - alpha) * pts1 + alpha * pts2

    # Add image corners and edges to cover the full frame
    h, w = img1.shape[:2]
    rect_pts = np.array([[0, 0], [w//2, 0], [w-1, 0],
                          [0, h//2], [w-1, h//2],
                          [0, h-1], [w//2, h-1], [w-1, h-1]],
                         dtype=np.float64)
    pts1 = np.vstack([pts1, rect_pts])
    pts2 = np.vstack([pts2, rect_pts])
    pts_mean = np.vstack([pts_mean, rect_pts])

    # Delaunay triangulation on mean shape
    triangles = get_delaunay_triangles(pts_mean)

    output = np.zeros_like(img1)
    for tri in triangles:
        src_tri1 = pts1[tri].astype(np.float32)
        src_tri2 = pts2[tri].astype(np.float32)
        dst_tri = pts_mean[tri].astype(np.float32)

        w1 = warp_triangle(img1, src_tri1, dst_tri, (w, h))
        w2 = warp_triangle(img2, src_tri2, dst_tri, (w, h))

        output = cv2.addWeighted(output, 1.0,
                                (1 - alpha) * w1 + alpha * w2, 1.0, 0)
    return output
\end{minted}

\subsubsection*{Generating a Morph Sequence}

\begin{minted}{python}
def generate_morph_sequence(img1, img2, lm1, lm2, num_frames=30,
                            output_dir="morph_frames"):
    """Generate a sequence of frames morphing from img1 to img2."""
    os.makedirs(output_dir, exist_ok=True)
    pts1 = lm1.astype(np.float64)
    pts2 = lm2.astype(np.float64)

    # Add rectangle points (same as morph_faces)
    h, w = img1.shape[:2]
    rect_pts = np.array([[0, 0], [w//2, 0], [w-1, 0],
                          [0, h//2], [w-1, h//2],
                          [0, h-1], [w//2, h-1], [w-1, h-1]],
                         dtype=np.float64)
    pts1 = np.vstack([pts1, rect_pts])
    pts2 = np.vstack([pts2, rect_pts])

    for k in range(num_frames):
        t = k / (num_frames - 1)
        frame = morph_faces(img1, img2, pts1, pts2, alpha=t)
        cv2.imwrite(f"{output_dir}/frame_{k:04d}.png", frame)
        print(f"  Frame {k+1}/{num_frames}")

    # Create video (requires ffmpeg)
    os.system(
        f"ffmpeg -framerate 30 -i {output_dir}/frame_%04d.png "
        f"-c:v libx264 -pix_fmt yuv420p -y morph.mp4"
    )
    print(f"Video saved as morph.mp4 in {output_dir}/")
\end{minted}

\subsubsection*{Command-Line Interface}

\begin{minted}{python}
def main():
    parser = argparse.ArgumentParser(
        description="Feature-based face morphing")
    parser.add_argument("image1", help="Path to first face image")
    parser.add_argument("image2", help="Path to second face image")
    parser.add_argument("--alpha", type=float, default=0.5,
        help="Blending weight (0=pure image1, 1=pure image2)")
    parser.add_argument("--backend", choices=["dlib", "mediapipe"],
        default="mediapipe",
        help="Landmark detection backend")
    parser.add_argument("--sequence", action="store_true",
        help="Generate morph sequence instead of single frame")
    parser.add_argument("--num-frames", type=int, default=30,
        help="Number of frames for sequence")
    parser.add_argument("--output", default="output.png",
        help="Output file path")
    parser.add_argument("--dlib-model", default="shape_predictor_68_face_landmarks.dat",
        help="Path to dlib shape predictor model (dlib backend only)")
    args = parser.parse_args()

    img1 = cv2.imread(args.image1)
    img2 = cv2.imread(args.image2)
    if img1 is None or img2 is None:
        print("Error: could not read input images")
        return

    if args.backend == "dlib":
        get_lm = lambda img: get_landmarks_dlib(img, args.dlib_model)
    else:
        get_lm = get_landmarks_mediapipe

    lm1 = get_lm(img1)
    lm2 = get_lm(img2)

    if args.sequence:
        generate_morph_sequence(img1, img2, lm1, lm2,
                                args.num_frames)
    else:
        result = morph_faces(img1, img2, lm1, lm2, args.alpha)
        assert not cv2.imwrite(args.output, result), \
            f"Failed to save output to {args.output}"
        print(f"Saved morph to {args.output}")

if __name__ == "__main__":
    main()
\end{minted}

\subsection{Usage Examples}

\begin{itemize}
    \item \textbf{Single frame (average face):}
    \begin{minted}{bash}
python face_morph.py face_A.jpg face_B.jpg \
    --backend mediapipe --alpha 0.5 --output average.png
    \end{minted}

    \item \textbf{Full morph sequence (30 frames):}
    \begin{minted}{bash}
python face_morph.py face_A.jpg face_B.jpg \
    --backend mediapipe --sequence --num-frames 30
    \end{minted}

    \item \textbf{Custom blend weight and dlib backend:}
    \begin{minted}{bash}
python face_morph.py face_A.jpg face_B.jpg \
    --backend dlib --dlib-model shape_predictor_68_face_landmarks.dat \
    --alpha 0.7 --output blend.png
    \end{minted}
\end{itemize}

\subsection{Complexity Analysis}

Let $N$ be the number of landmark points (68 for dlib, 68 mapped from MediaPipe), $W \times H$ be the image dimensions, and $T \approx 2N$ be the number of Delaunay triangles.

\begin{table}[h]
\centering
\begin{tabularx}{\textwidth}{lccc}
\toprule
\textbf{Step} & \textbf{Complexity} & \textbf{Notes} \\
\midrule
Landmark detection & $O(WH)$ & CNN inference (dominant for large images) \\
Delaunay triangulation & $O(N \log N)$ & Using Bowyer-Watson with spatial index \\
Affine warp computation & $O(N)$ & 3x3 matrix inversion per triangle \\
Inverse mapping + sampling & $O(WH \cdot N)$ & Naive: each pixel searches all triangles; \\
& & with spatial indexing: $O(WH)$ \\
Blending & $O(WH)$ & Per-pixel alpha blend \\
\midrule
\textbf{Total (single frame)} & \textbf{$O(WH \cdot N)$} & & \\
\midrule
\textbf{Total (sequence, $F$ frames)} & \textbf{$O(F \cdot WH \cdot N)$} & & \\
\bottomrule
\end{tabularx}
\end{table}

In practice, the bounding box optimization reduces the effective complexity to $O(WH)$ because each triangle only processes pixels within its bounding box, and the point-in-triangle test is only performed for pixels that belong to that triangle. With OpenCV's GPU-accelerated \texttt{warpAffine}, a single frame at 1024$\times$1024 resolution takes approximately 0.1--0.3 seconds on a modern CPU and under 0.05 seconds on GPU.

\subsection{Common Pitfalls and Debugging Tips}

\begin{enumerate}
    \item \textbf{Face detector misses the face.} This is the most common failure. Ensure faces are roughly front-facing, well-lit, and at least 200$\times$200 pixels. If dlib fails, try the MediaPipe backend which has better multi-pose detection.

    \item \textbf{Triangulation produces triangles outside the face.} The Delaunay triangulation of the landmarks alone will include triangles in the background (e.g., forehead, neck, sides of head). This is why we add the eight rectangle points (corners and edge midpoints) to the landmark set. These ``anchor'' points push the triangulation outward so that the convex hull covers the entire image.

    \item \textbf{Visible seams between triangles.} Seams appear when the warped colors from adjacent triangles don't match at their shared boundary. This is typically caused by:
    \begin{itemize}
        \item Landmark detection errors (points are slightly off anatomical positions)
        \item Rasterization gaps (rounding errors in triangle filling)
    \end{itemize}
    A 1-pixel Gaussian blur on the output often eliminates visible seams.

    \item \textbf{Holes in the output.} Holes indicate that forward mapping is being used instead of inverse mapping. Verify that you are using \texttt{warpAffine} with the \emph{inverse} matrix (mapping destination $\to$ source). OpenCV's \texttt{warpAffine} expects the forward matrix (source $\to$ destination), but since we need inverse mapping, we compute \texttt{inv\_warp = compute\_affine\_transform(dst\_tri, src\_tri)}.

    \item \textbf{Color mismatch after blending.} If the two input faces have different lighting, the blended result may look unnatural. Apply color correction before blending: compute the mean and standard deviation of each RGB channel in both faces, then adjust the brighter face to match the darker one.

    \item \textbf{Slow performance.} For morph sequences, the bottleneck is Delaunay recomputation at every frame. If you need real-time performance:
    \begin{itemize}
        \item Pre-compute the triangulation on the mean shape and reuse it (triangles stay fixed)
        \item Use \texttt{cv2.Subdiv2D} instead of SciPy for faster triangulation
        \item Downscale input images, morph, then upscale the result
    \end{itemize}

    \item \textbf{MediaPipe 68-point mapping inaccuracy.} The MediaPipe-to-68-point mapping is an approximation. For the highest quality results with MediaPipe, use all 478 landmarks directly and compute the triangulation on the full set. This produces a much denser mesh (thousands of triangles vs. $\sim$130) and smoother warps.
\end{enumerate}

\subsection{Summary}

The feature-based face morphing pipeline consists of five interconnected steps:
\begin{enumerate}
    \item \textbf{Detect} facial landmarks on both input faces
    \item \textbf{Interpolate} the landmark positions to the desired mean shape
    \item \textbf{Triangulate} using Delaunay to decompose the face into regions
    \item \textbf{Warp} each face to the mean shape using piecewise affine transforms
    \item \textbf{Blend} the warped faces with alpha compositing
\end{enumerate}

Each step builds on mathematical foundations — affine geometry, barycentric coordinates, and computational geometry — that we have derived from first principles in the preceding chapters. The implementation presented here is complete and ready to use, supporting both dlib and MediaPipe landmark backends through a unified command-line interface.

For further reading, the foundational papers are:
\begin{itemize}
    \item Beier and Neely, ``Feature-Based Image Metamorphosis,'' SIGGRAPH 1992 --- the original paper that introduced this approach
    \item Wolberg, ``Image Morphing: A Survey,'' The Visual Computer, 1998 --- comprehensive survey of morphing techniques
    \item Roweis and Saul, ``Nonlinear Dimensionality Reduction by Locally Linear Embedding,'' Science 2000 --- related techniques for manifold-based face representation
\end{itemize}
